% ----------------------
\subsection{Section 93 (6 May 1833)}
% ----------------------
	\label{DC93}
	
	% -----
	\subsubsection{Historical Background}
	% -----
		\label{DC93-HB}
		
		% --
		\paragraph{JSP}
		% --
		
			``Few contemporary sources shed light on the background of this 6 May 1833 revelation. JS's history simply states, `On the 6\supersu{th.} [of May 1833] I received the following.' Newel K. Whitney, one of the individuals mentioned in the text, made the only other known early comment on the back of perhaps the earliest extant manuscript version of the revelation. Whitney summarized the revelation by writing, `Revelation to Joseph, Sidny Frederick \& Newell by chastisement \& also relative to the Father \& Son 6 May 1833.' While Whitney's summary focused first on the reprimand the men received from the Lord in relation to their families, the majority of the revelation was `relative to the Father \& Son.'

			``The text of the revelation appears to be closely related to the first chapter of the Gospel according to John but was likely not the direct result of JS's work revising the New Testament, since the revision had been completed three months earlier, on 2 February 1833. Further, the revisions to John's record made in this revelation generally vary from the changes JS made to the same text in his larger project of revising the Bible. The revelation was not the first JS document to involve passages from John's gospel. One of the earliest recorded JS revelations, dictated April 1829, was a result of a disagreement between JS and Oliver Cowdery over the final chapter of John and whether Jesus Christ's statement about John tarrying meant that he would not die. Through the means of a seer stone or stones in his possession, JS dictated the 1829 revelation, which he explained had been `translated from parchment' and included a fuller account from John, which John had `written and hid up.' The Book of Mormon also references John by name several times, while no other apostle of Jesus is specifically mentioned in the book. In addition, in March 1832 JS created a document that explained the content of many verses in the book of Revelation, which the Book of Mormon explicitly asserts was authored by John the Apostle.

			``The revelation featured here directly challenged several prevailing Christian beliefs of the time, including doctrines regarding the nature of Jesus Christ, especially his humanity and divinity, that most Christians believed had been settled by the Council of Chalcedon in AD 451. That council held that Jesus Christ was both fully human and fully divine, `that in Christ two distinct natures were united in one person, without any change, mixture, or confusion.' This revelation instead describes Jesus as having `received not of the fulness at the first but received grace for grace and he received not of the fulness but continued from grace to grace until he received a fulness.'

			``Aside from addressing the nature of Jesus, the revelation also addressed humankind's relationship with God, asserting that `man was also in the begining with God.' The premortal existence of humankind was also discussed in an earlier document dictated by JS; in his early translation of the Bible, he recorded a conversation between God and Moses in which Moses was told, `For I the Lord God created all things of which I have spoken spiritually before they were naturally upon the face of the Earth . . . \& I the Lord God had created all the children of men \& not yet a man to till the ground for in Heaven created I them \& there was not yet flesh upon the Earth.' The 6 May revelation also challenged the general Christian concept of creation ex nihilo by stating that `inteligence or the Light of truth was not created or made.' In an even greater contrast to the prevailing belief in the omnipotence of God, this revelation expressly stated not only that intelligence was `not created or made' by God, but also that it could not, in fact, be created by him.

			``It is not known whether any discussions about early Christian creeds and doctrines preceded this revelation or how familiar JS was with them prior to dictating it. Sidney Rigdon, one of the revelation's recipients, had been a follower of Alexander Campbell's religious movement that denied that belief in any of the traditional Christian creeds was essential. Campbell had also published a discourse in 1827, centered on the first chapter of John, rejecting the traditional Calvinist interpretation. In addition, the Mormon settlement in Kirtland, Ohio, was not far from a large Shaker community in North Union, Ohio. The Shakers' belief in the nature of Christ, widely viewed as heretical by orthodox Christians of the time, had in part been the subject of an 1831 revelation that commanded missionaries to preach to the Shakers. These influences, as well as many others, may have led church members to discuss topics such as the nature of God and Jesus before this revelation was dictated.

			``This revelation was first published in the 1835 edition of the Doctrine and Covenants. By 1839, JS appears to have expanded upon some of the teachings it contained. In a sermon that year, he reportedly said, `The spirit of man is not a created being: it existed from Eternity and will exist to Eternity. Any thing created cannot be eternal and Earth, water, \&c all these had their existence in an elementary state from Eternity. Our Savior speaks of children and says their Angels always stand before my Father. The Father called all spirits before him at the creation of man and organized them.'%
				\footnote{%
					% \label{}%
					See \hyperref[EMS-1839-JS-CIR-26-JUN-CIR-4-AUG-1839-A]{this discourse}.
					}
			A few years later he further explained, `God was a self exhisting being, man exhist upon the same principle. God made a tabernacle \& por spirit in it and it became a Human soul, man exhisted in spirit \& mind coequal with God himself. . . . God has power to institute laws to instruct the weaker intelligences that they may be exhalted with himself.'%
				\footnote{%
					% \label{}%
					This is a reference to WW's account of the King Follett sermon on 7 April 1844.
					}

			``Although three early manuscript versions of this revelation are extant, the one featured here, from the Newel K. Whitney collection, is most likely the earliest version. As the revelation was directed in part to Whitney, he likely received this copy immediately after it was dictated by JS. It is even possible that Whitney's version is the original copy, recorded while JS dictated the revelation. In any case, the Whitney version is very closely related to the copy in Revelation Book 2, even including similar scribal errors and omissions. Orson Hyde recorded the copy in Revelation Book 2 as early as 1 June, though he most likely recorded the revelation after his appointment as scribe on 6 June 1833. Hyde apparently used Whitney's copy of the revelation to create the Revelation Book 2 version. The copy found in Revelation Book 1 was almost certainly copied later than these two versions because that revelation book was still in Missouri with John Whitmer at this time. In a few places the version in Revelation Book 1 includes words and phrases not found in the other two versions, possibly reflecting some editorial changes made by Whitmer when he copied it into the book. In the following transcript [in the JSP], significant differences between versions in the revelation books and the version featured here are noted.''
			
	% -----
	\subsubsection{Versions}
	% -----
		\label{DC93-V}
		
		See the \href{https://drive.google.com/file/d/1goAvNz8jS4R8slYfMG_db55Ty2z_vmgK/view?usp=sharing}{sections comparison} document.
	
		\begin{compactdesc}
			\item[JSP] \hyperref[EMS-1833-JS-6-MAY-1833]{6 May 1833}
			\item[BC] ---
			\item[DC 1835] Section 82, 210--213
			\item[DC 1844] Section 83, 323--327
			\item[TS] 6:1 (15 January 1845), 768--769
			\item[MS] 14:27 (28 August 1852), 422--424
		\end{compactdesc}

	% -----
	\subsubsection{Section 93:1} % *DC93-1 
	% -----
		\label{DC93-1}

			\footnote{%
				% \label{}%
				For a discussion of the structure of this chapter, see my \hyperref[EPIST-REVELATION-DC]{notes on epistemology for the DC}.
				}%
		VERILY, thus saith the Lord:
			\footnote{%
				% \label{}%
				There are five requirements given here to realize the promise of seeing Christ's (?) face and knowing that he is:
					\begin{compactenum}
						\item Forsake your sins
						\item Come unto Christ
						\item Call on Christ's name
						\item Obey his voice
						\item Keep (all [see JSP]) Christ's commandments
					\end{compactenum}
				That's a lot of stuff to do over a long period of time, but I see this is as the goal of righteousness---you are working toward ever greater manifestations because of your righteousness.
				}%
		It shall come to pass that every soul who forsaketh his sins and cometh unto me, and calleth on my name, and obeyeth my voice, and keepeth%
			\footnote{%
				% \label{}%
				The JSP version has ``all'' here: ``keepeth \textit{all} my commandments.''
				}
		my commandments, shall see my face and know that I am;

		\begin{framed}
		\scriptsize
			\begin{compactdesc}
				\item[JSP] Verely thus saith the Lord, it shall come to pass, that evry soul who forsaketh their sins and cometh unto me and calleth on my name and obeyeth my voice and keepeth all my commandments shall see my face and know that I am
				\item[BC] ---
				\item[RB2] Verily thus saith the Lord, it shall come to pass that every soul who forsaketh their sins and cometh unto me and calleth on my name and obeyeth my voice and keepeth all my commandments, shall see my face and know that I am, 
				\item[RB1] Verily, thus saith the Lord, it shall come to pass, that every soul who forsaketh their sins, \& cometh unto me, \& calleth on my name, \& obeyeth my voice, \& keepeth all my commandments, shall see my face \& know that I am, 
				\item[DC 1835] 1 Verily thus saith the Lord, it shall come to pass that every soul who forsaketh their sins and cometh unto me, and calleth on my name, and obeyeth my voice, and keepeth my commandments, shall see my face, and know that I am, 
				% \item[DC 1844]
			\end{compactdesc}
		\end{framed}

	% -----
	\subsubsection{Section 93:2} % *DC93-2 
	% -----
		\label{DC93-2}

		And that I am the true light that lighteth every man%
		\footnote{%
			% \label{}%
			Compare \hyperref[DC88-5]{DC 88:5--13}, especially verse 11.
			} 
		that cometh into the world;

		\begin{framed}
		\scriptsize
			\begin{compactdesc}
				\item[JSP] and that I am the true light that lighteth evry man who cometh into the world;
				\item[BC] ---
				\item[RB2] and that I am the true light that lighteth every man who cometh into the world, 
				\item[RB1] \sout{\& that} <\& that> I am the true light that lighteneth every man who cometh into the world 
				\item[DC 1835] and that I am the true light that lighteth every man that cometh into the world: 
				% \item[DC 1844] 
				% \item[]
			\end{compactdesc}
		\end{framed}

	% -----
	\subsubsection{Section 93:3} % *DC93-3 
	% -----
		\label{DC93-3}

		And that I am in the Father, and the Father in me, and the Father and I are one---

		\begin{framed}
		\scriptsize
			\begin{compactdesc}
				\item[JSP] and that I am in the fathe[r] and the father in me and the fathe[r] and I are one
				\item[BC] ---
				\item[RB2] and that I am in the Father and the Father in me, and the Father and I are one.
				\item[RB1] \& that I am <in> the Father, \sout{because he gave} \& the Father in me: And the Father \& I are one. 
				\item[DC 1835] and that I am in the Father and the Father in me, and the Father and I are one: 
				% \item[DC 1844] 
				% \item[]
			\end{compactdesc}
		\end{framed}

	% -----
	\subsubsection{Section 93:4} % *DC93-4 
	% -----
		\label{DC93-4}

		The Father because he gave me of his \hyperref[FULNESS]{fulness}, and the Son%
			\footnote{%
				% \label{}%
				Compare \hyperref[MOSIAH-15-2]{Mosiah 15:2--5}. 
				}
		because I was in the world and made flesh my tabernacle, and dwelt among the sons of men.

		\begin{framed}
		\scriptsize
			\begin{compactdesc}
				\item[JSP] the father because he gave me of his fulness and the son becaus I was in the world and made flesh my tabernacle and dwelt among the sons of men
				\item[BC] ---
				\item[RB2] The Father because he gave me of his fullness, and the son because I was in the world and made flesh my tabernacle and dwelt among the sons of men. 
				\item[RB1] The Father, because he gave me of his fulness; and the son, because I was in the world, \& made flesh my tabernacle, \& dwelt among the Sons of men. 
				\item[DC 1835] the Father because he gave me of his fulness; and the Son because I was in the world and made flesh my tabernacle, and dwelt among the sons of men: 
				% \item[DC 1844] 
				% \item[]
			\end{compactdesc}
		\end{framed}

	% -----
	\subsubsection{Section 93:5} % *DC93-5 
	% -----
		\label{DC93-5}

		I was in the world and received of my Father, and the works of him were plainly manifest.

		\begin{framed}
		\scriptsize
			\begin{compactdesc}
				\item[JSP] I was in the world and received of my father, and the works of him were plainly manifest
				\item[BC] ---
				\item[RB2] I was in the world and received of my Father, and the works of him were plainly manifest
				\item[RB1] I was in the world, and received of my Father, \& <the> works of him were plainly manifest: 
				\item[DC 1835] I was in the world and received of my Father, and the works of him were plainly manifest; 
				% \item[DC 1844] 
				% \item[]
			\end{compactdesc}
		\end{framed}

	% -----
	\subsubsection{Section 93:6} % *DC93-6 
	% -----
		\label{DC93-6}

		And John saw and bore record of the \hyperref[FULNESS]{fulness} of my glory, and the \hyperref[FULNESS]{fulness} of John's record is \hyperref[HEREAFTER]{hereafter} to be revealed.%
			\footnote{%
				% \label{}%
				Does this last part mean that the fulness of John's record is going to be revealed at some point later on? Or does it mean that what follows, starting in verse 7, \textit{is} the fulness of John's record? Definitely could be the latter. The word ``hereafter'' makes it seem like it's farther in the future though.
				}

		\begin{framed}
		\scriptsize
			\begin{compactdesc}
				\item[JSP] and John saw and bear record of the fulness of \sout{his} <my> glory and the fulness of Johns reccord is hereafter to be reveiled
				\item[BC] ---
				\item[RB2] and I John saw and bare reccord of the fullness of my glory, and the fullness of John's reccord is here after to be revealed, 
				\item[RB1] And John saw \& bore record of the fulness of my glory. And the fulness of \sout{my glory} John's record is hereafter to be revealed. 
				\item[DC 1835] and John saw and bore record of the fulness of my glory; and the fulness of John’s record is hereafter to be revealed: 
				% \item[DC 1844] 
				% \item[]
			\end{compactdesc}
		\end{framed}

	% -----
	\subsubsection{Section 93:7} % *DC93-7 
	% -----
		\label{DC93-7}

		And he bore record, saying: I saw his glory, that he was in the beginning, before the world was;%
			\footnote{%
				% \label{}%
				The punctuation, which is absent in the JSP version, has a really strong influence on how we read this verse. For example, consider how we'd read it if it were punctuated like this: ``I saw his glory that he was, in the beginning, before the world was.''
				}

		\begin{framed}
		\scriptsize
			\begin{compactdesc}
				\item[JSP] and he bear reccord saying I saw his glory that he was in the begining before the world was
				\item[BC] ---
				\item[RB2] and he bare reccord saying, ``I saw his glory that he was in the beginning before the world was. 
				\item[RB1] And he bore record saying, I saw his glory, that he was in the beginning, before the world was. 
				\item[DC 1835] And he bore record, saying, I saw his glory that he was in the beginning before the world was: 
				% \item[DC 1844] 
				% \item[]
			\end{compactdesc}
		\end{framed}

	% -----
	\subsubsection{Section 93:8} % *DC93-8 
	% -----
		\label{DC93-8}

		Therefore, in the beginning the Word was, for he was the Word, even the
			\footnote{%
			% \label{}%
			The title ``messenger of salvation'' doesn't appear anywhere else in the scriptures.
			}%
		messenger%
			\footnote{%
				% \label{}%
				\gr{ἄγγελος}, in Greek; BDAG \#7.a, ``\textbf{a transcendent power who carries out various missions or tasks, \textit{messenger, angel}}.'' OED \#1.a, ``A person who carries a message or goes on an errand for another; a courier. Formerly also: an envoy or ambassador; a person who brings news or other intelligence; a spy, a scout (\textit{obsolete}).'' I think the ``for another'' is key to understanding this title or role as ``messenger of salvation.'' Christ is the one sent here to carry out the work of salvation for or in behalf of the Father; he is the one we look to for our salvation, the representative of salvation on earth. Easy to connect this to the priesthood as well.
				}
		of salvation---,

		\begin{framed}
		\scriptsize
			\begin{compactdesc}
				\item[JSP] therefore in the beg[inn]ing the word was for he was the word even the messenger of salvation
				\item[BC] ---
				\item[RB2] Therefore in the beginning the word was, for He was the word even the messenger of salvation, 
				\item[RB1] Therefore, in the beginning the word was; for he was the word, even the messenger of Salvation; 
				\item[DC 1835] therefore, in the beginning the Word was; for he was the Word, even the messenger of salvation, 
				% \item[DC 1844] 
				% \item[]
			\end{compactdesc}
		\end{framed}

	% -----
	\subsubsection{Section 93:9} % *DC93-9 
	% -----
		\label{DC93-9}

		The light and the Redeemer of the world; the Spirit of truth,%
			\footnote{%
				% \label{}%
				For ``spirit of truth,'' see \hyperref[SPIRIT-PHRASES]{here}.
				}
		who came into the world, because the world was made by him, and in him was the life of men and the light of men.

		\begin{framed}
		\scriptsize
			\begin{compactdesc}
				\item[JSP] the light and the redeemer of the world the spirit of truth who came into the world becaus the world was made by him and in him was the \sout{<fife>} \sout{light} <Life> of men and the light of men
				\item[BC] ---
				\item[RB2] the light and the redeemer of the world; the spirit of truth who came into the world because the world was made by him, and in him was the life of men, and the light of men. 
				\item[RB1] the light \& the Redeemer of the world, the Spirit of truth who came into the world, because the world was made by him. And in him was the life of men \& the light of men. 
				\item[DC 1835] the light and the Redeemer of the world; the Spirit of truth, who came into the world because the world was made by him; and in him was the life of men and the light of men. 
				% \item[DC 1844] 
				% \item[]
			\end{compactdesc}
		\end{framed}

	% -----
	\subsubsection{Section 93:10} % *DC93-10 
	% -----
		\label{DC93-10}

			\footnote{%
				% \label{}%
				Compare this verse with \hyperref[DC76-24]{DC 76:24}.
				}%
		The worlds were made by him; men were made by him; all things were made by him, and through him, and of him.

		\begin{framed}
		\scriptsize
			\begin{compactdesc}
				\item[JSP] the worlds were made by him men were made by him all things were made by him and through him and of him
				\item[BC] ---
				\item[RB2] The worlds were made by him Men were made by him; All things were made by him and through him and of him.
				\item[RB1] The worlds were made by him: Men were made by him: All things were made by him, \& through him \& of him. 
				\item[DC 1835] The worlds were made by him. Men were made by him. All things were made by him, and through him, and of him: 
				% \item[DC 1844] 
				% \item[]
			\end{compactdesc}
		\end{framed}

	% -----
	\subsubsection{Section 93:11} % *DC93-11 
	% -----
		\label{DC93-11}

		And I, John, bear record that I beheld his glory, as the glory of the Only Begotten of the Father, full of grace and truth, even the \hyperref[SPIRIT-PHRASES]{Spirit of truth}, which came and dwelt in the flesh, and dwelt among us.

		\begin{framed}
		\scriptsize
			\begin{compactdesc}
				\item[JSP] and I, John bear reccord that I beheld his glory as the glory of the only begotten of the fathe[r] full of grace and truth even the spirit of truth which came and dwelt in flesh and dwelt among us
				\item[BC] ---
				\item[RB2] and I John bear reccord that I beheld his glory, as the glory of the onely begotten of the father full of grace and truth, even the spirit of truth which came and dwelt in flesh, and dwelt among us. 
				\item[RB1] And I John bear record that I beheld his glory, as the glory of the only Redeemer. Begotten of the Father, full of grace \& truth; Even the Spirit of truth which came \& dwelt in flesh; \& dwelt among us. 
				\item[DC 1835] And I John bare record that I beheld his glory, as the glory of the Only Begotten of the Father, full of grace and truth; even the Spirit of truth which came and dwelt in the flesh, and dwelt among us.
				% \item[DC 1844] 
				% \item[]
			\end{compactdesc}
		\end{framed}

	% -----
	\subsubsection{Section 93:12} % *DC93-12 
	% -----
		\label{DC93-12}

		And I, John, saw that he received not of the \hyperref[FULNESS]{fulness} at the first, but received
			\footnote{%
				% \label{}%
				For ``grace for grace,'' see \hyperref[JOHN-1-16]{John 1:16} and notes there.
				}%
		grace for grace;%
			\footnote{%
				% \label{}%
				See the note above about ``grace for grace.'' There are a couple other verses which mention it in this section, but the place you want to compare here is \hyperref[JOHN-1-16]{John 1:16}. The simplest interpretation of this verse, and I think the correct one, is that Christ did not always possess a fulness of power and glory and grace---he received it a little at a time. This is very important for thinking about eternal progression, the priesthood, the nature of God, and many other topics. See also \hyperref[LUKE-2-40]{Luke 2:40} and \hyperref[LUKE-2-52]{Luke 2:52}, which notes that a person can ``increase'' \gr{χάριτι παρὰ θεῷ καὶ ἀνθρώποις}. That's what we need to do if we're going to get this ``grace'' or favor from God to change our nature over time---we have to increase in \gr{χάρις} with both God and people, and the way to do this is through living a life full of \gr{δικαιοσύνη}.
				} 

		\begin{framed}
		\scriptsize
			\begin{compactdesc}
				\item[JSP] and I John saw that he received not of the fulness at the first but received grace for grace
				\item[BC] ---
				\item[RB2] and I John saw that he received not of the fullness at the first, but received grace for grace, 
				\item[RB1] And I John saw that he received not of the fulness at <the> first; but received grace for grace. 
				\item[DC 1835] 2 And I John saw that he received not of the fulness at the first, but received grace for grace: 
				% \item[DC 1844] 
				% \item[]
			\end{compactdesc}
		\end{framed}

	% -----
	\subsubsection{Section 93:13} % *DC93-13 
	% -----
		\label{DC93-13}

		And he received not of the \hyperref[FULNESS]{fulness} at first, but continued from grace to grace,%
			\footnote{%
				% \label{}%
				For ``grace to grace,'' see \hyperref[JOHN-1-16]{John 1:16} and notes there.
				}
		until he received a \hyperref[FULNESS]{fulness};%
			\footnote{%
				% \label{}%
				What does this verse add to our understanding that isn't already there from the previous verse? 
				}

		\begin{framed}
		\scriptsize
			\begin{compactdesc}
				\item[JSP] and he received not of the fulness but continued from grace to grace until he received a fulness
				\item[BC] ---
				\item[RB2] and he received not of the fullness but continued from grace to grace untill he received a fullness, 
				\item[RB1] And he received not of the fulness; but continued from grace to grace untill he received a fulness.
				\item[DC 1835] and he received not of the fulness at first, but continued from grace to grace, until he received a fulness: 
				% \item[DC 1844] 
				% \item[]
			\end{compactdesc}
		\end{framed}

	% -----
	\subsubsection{Section 93:14} % *DC93-14 
	% -----
		\label{DC93-14}

		And thus he was called the Son of God, because he received not of the \hyperref[FULNESS]{fulness} at the first.%
			\footnote{%
				% \label{}%
				This is probably the clearest explanation in the scriptures of why Christ is known as ``the Son.'' Note that the explanation does \textit{not} involve some kind of parent-child relationship.
				}

		\begin{framed}
		\scriptsize
			\begin{compactdesc}
				\item[JSP] and thus he was called the son of God because he received not of the fulness at the first
				\item[BC] ---
				\item[RB2] and thus he was called the son of God because he received not of the fullness at first. 
				\item[RB1] [paragraph break] And thus he was called the Son of God, because he received not of the fulness at the first. 
				\item[DC 1835] and thus he was called the Son of God, because he received not of the fulness at the first. 
				% \item[DC 1844] 
				% \item[]
			\end{compactdesc}
		\end{framed}

	% -----
	\subsubsection{Section 93:15} % *DC93-15 
	% -----
		\label{DC93-15}

		And I, John, bear record, and lo, the heavens were opened, and the Holy Ghost descended upon him in the form of a dove, and sat upon him, and there came a voice out of heaven saying: This is my beloved Son.

		\begin{framed}
		\scriptsize
			\begin{compactdesc}
				\item[JSP] and I John bear reccord and lo the heavens were opened and the holy ghost decended upon him in the form of a dove and set upon him and there came a voice out of heaven saying this is my beloved son,
				\item[BC] ---
				\item[RB2] and I John bear reccord and lo: the heavens were opened and the holy ghost descended upon him in the form of a dove and set upon him, and there came a voice out of heaven saying, this is my beloved son
				\item[RB1] And I John bear record \& lo, the Heavens were opened, and the Holy Ghost decended upon him in the form of a dove, \& Set upon him. And there came a voice out of heaven saying, this is my beloved son. 
				\item[DC 1835] And I John bare record, and lo, the heavens were opened and the Holy Ghost descended upon him in the form of a dove, and sat upon him, and there came a voice out of heaven saying, this is my beloved Son. 
				% \item[DC 1844] 
				% \item[]
			\end{compactdesc}
		\end{framed}

	% -----
	\subsubsection{Section 93:16} % *DC93-16 
	% -----
		\label{DC93-16}

		And I, John, bear record that he received a \hyperref[FULNESS]{fulness} of the glory of the Father;

		\begin{framed}
		\scriptsize
			\begin{compactdesc}
				\item[JSP] and I John bear reccord that he received a fulness of the glory of the \sout{eternal} \sout{God} father
				\item[BC] ---
				\item[RB2] and I John bare reccord that he received a fullness of the glory of the Father
				\item[RB1] And I John, bear record, that he received of fulness of the glory of the Father. 
				\item[DC 1835] And I John bare record that he received a fulness of the glory of the Father; 
				% \item[DC 1844] 
				% \item[]
			\end{compactdesc}
		\end{framed}

	% -----
	\subsubsection{Section 93:17} % *DC93-17 
	% -----
		\label{DC93-17}

		And he received all \hyperref[POWER]{power},%
			\footnote{%
				% \label{}%
				In the verse in Matthew, it's \gr{ἐξουσία}.
				}
		both in heaven and on earth,%
			\footnote{%
				% \label{}%
				See \hyperref[MATTHEW-28-18]{Matthew 28:18}.
				}
		and the glory of the Father was with him, for he dwelt in him.%
			\footnote{%
				% \label{}%
				A clear statement of the in-dwelling relationship.
				}

		\begin{framed}
		\scriptsize
			\begin{compactdesc}
				\item[JSP] and he receivd all power both in heaven and on earth and the glory of the father was with him for he dwelt in him
				\item[BC] ---
				\item[RB2] and he received all power both in heaven and on earth and the glory of the father was with him, for he dwelt in him
				\item[RB1] And he received all power, both in heaven and on the earth. And the glory of the Father was with him for he dwelt in him. 
				\item[DC 1835] and he received all power, both in heaven and on earth; and the glory of the Father was with him, for he dwelt in him.
				% \item[DC 1844] 
				% \item[]
			\end{compactdesc}
		\end{framed}

	% -----
	\subsubsection{Section 93:18} % *DC93-18 
	% -----
		\label{DC93-18}

		And it shall come to pass, that if you are faithful you shall receive the \hyperref[FULNESS]{fulness} of the record of John.

		\begin{framed}
		\scriptsize
			\begin{compactdesc}
				\item[JSP] and it shall come to pass that if you are faithful you shall receive the fulness of the reccord of John
				\item[BC] ---
				\item[RB2] and it shall come to pass that if you are faithful you shall receive the fullness of the reccord of John,
				\item[RB1] And it shall come to pass, that if you are faithful, you shall receive the fulness of the record of John. 
				\item[DC 1835] 3 And it shall come to pass that if you are faithful, you shall receive the fulness of the record of John. 
				% \item[DC 1844] 
				% \item[]
			\end{compactdesc}
		\end{framed}

	% -----
	\subsubsection{Section 93:19} % *DC93-19 
	% -----
		\label{DC93-19}

			\footnote{%
				% \label{}%
				An incredibly important verse that tells you why Christ is revealing everything that has come before.
				}%
		I give unto you these sayings that you may understand and know how to worship,%
			\footnote{%
				% \label{DC93-19-how}%
				See \hyperref[REPENTANCE-learning]{this study}.
				}
		and know what you worship,%
			\footnote{%
				% \label{DC93-19-what}%
				See \hyperref[JOHN-4-22]{John 4:22}.
				}
		that you may come unto the Father in my name, and in due time receive of his \hyperref[FULNESS]{fulness}.

		\begin{framed}
		\scriptsize
			\begin{compactdesc}
				\item[JSP] I give unto you these sayings that you may understand and know how to worship and know what you worship that you may come unto the fathe[r] in my name and in due time receive of his fulness
				\item[BC] ---
				\item[RB2] I give unto you these sayings that you may understand and know how to worship, and know what you worship that you may come unto the father in my name, and in due time receive of his fullness,
				\item[RB1] I give unto you these sayings, that you may understand, and know how to worship, that you may come unto the Father in my name, and in due time receive of his fulness. 
				\item[DC 1835] I give unto you these sayings that you may understand and know how to worship, and know what you worship, that you may come unto the Father in my name, and in due time receive of his fulness 
				% \item[DC 1844] 
				% \item[]
			\end{compactdesc}
		\end{framed}

	% -----
	\subsubsection{Section 93:20} % *DC93-20 
	% -----
		\label{DC93-20}

		For if you keep my commandments you shall receive of his fulness, and be glorified in me as I am in the Father; therefore, I say unto you, you shall receive grace for grace.%
			\footnote{%
				% \label{}%
				So in other words, we will progress in the same way Christ does.
				}

		\begin{framed}
		\scriptsize
			\begin{compactdesc}
				\item[JSP] for if you keep my commandments you shall receive of his fulness and be glor[i]fied in me as I am glor[i]fied in the father, therefore I say <unto> you you shall receive grace for grace
				\item[BC] ---
				\item[RB2] for if you keep my commandments you shall receive of his fullness, and be glorified in me as I am glorified in the Father. Therefore, I say unto you, you shall receive grace for grace.
				\item[RB1] For if you keep my commandments you shall receive of his fulness, and be glorified in me, as I am glorified in the Father. Therefore, I say unto you, you shall receive grace for grace. 
				\item[DC 1835] for if you keep my commandments you shall receive of his fulness and be glorified in me as I am in the Father: therefore, I say unto you, you shall receive grace for grace. 
				% \item[DC 1844] 
				% \item[]
			\end{compactdesc}
		\end{framed}

	% -----
	\subsubsection{Section 93:21} % *DC93-21 
	% -----
		\label{DC93-21}

		And now, verily I say unto you, I was in the beginning with the Father, and am the \hyperref[FIRSTBORN]{Firstborn};

		\begin{framed}
		\scriptsize
			\begin{compactdesc}
				\item[JSP] and now verely I say unto you I was in the begining with the fathe[r] and am the first born
				\item[BC] ---
				\item[RB2]  And now verily I say unto you I was in the beginning with the father, and am the first born 
				\item[RB1] And now verily I say unto you, I was in the beginning with the father, and am the first born: 
				% \item[DC 1835] 
				% \item[DC 1844] 
				% \item[]
			\end{compactdesc}
		\end{framed}

	% -----
	\subsubsection{Section 93:22} % *DC93-22 
	% -----
		\label{DC93-22}

		And all those who are begotten through me are partakers of the glory of the same, and are the church of the \hyperref[FIRSTBORN]{Firstborn}.%
			\footnote{%
				% \label{}%
				For ``church of the Firstborn,'' see \hyperref[HEBREWS-12-23]{Hebrews 12:23} and \hyperref[CHURCH-PHRASES]{this phrase concordance}.
				}

		\begin{framed}
		\scriptsize
			\begin{compactdesc}
				\item[JSP] and all those who are begotten through me are partakers of the glory of the same and are the church of the first born,
				\item[BC] ---
				\item[RB2] and all those who are begotten through me are partakers of the glory of the same. and are the church of the first born 
				\item[RB1] and all those who are begotten throug[h] me, are partakers of the glory of the same, and are the church of the first born. 
				% \item[DC 1835] 
				% \item[DC 1844] 
				% \item[]
			\end{compactdesc}
		\end{framed}

	% -----
	\subsubsection{Section 93:23} % *DC93-23 
	% -----
		\label{DC93-23}

		Ye were also in the beginning with the Father; that which is Spirit, even the Spirit of truth;

		\begin{framed}
		\scriptsize
			\begin{compactdesc}
				\item[JSP] ye were also in the begining with the fathe[r] that which is Spirit even the spirit of truth
				\item[BC] ---
				\item[RB2] Ye were also in the beginning with the father, that which is spirit even the spirit of truth 
				\item[RB1] Ye were also in the beginning with the Father, that which is spirit, even the spirit of truth. 
				% \item[DC 1835] 
				% \item[DC 1844] 
				% \item[]
			\end{compactdesc}
		\end{framed}

	% -----
	\subsubsection{Section 93:24} % *DC93-24 
	% -----
		\label{DC93-24}

			\footnote{%
				% \label{}%
				This verse deserves a whole series of studies on its own, but maybe the most important aspect of it is that it subordinates truth to knowledge, where normally you think of the priority relationship as going in reverse; that's extremely important, including for the epistemology of revelation. It makes me think of Kant's first \textit{Critique} and its importance in thinking about the gospel. EDIT: And now, years after I originally wrote this footnote, I have written this series of studies, or at least the first preliminary series. Find them \href{https://docs.google.com/document/d/1Skh5HA1vFgllloVAObrsPw6zufYDzGo_C-vdkToHbGw/edit\#heading=h.c59zktk9yl84}{here}.
				}%
		And truth is
			\footnote{%
				% \label{}%
				See \hyperref[1NEPHI-1-3]{1 Nephi 1:3}, 
				}%
		knowledge of things as they are, and as they were, and as they are to come;%
			\footnote{%
				% \label{}%
				See \hyperref[REVELATIONNT-1-4]{Revelation 1:4} and notes there. That verse speaks of God as ``him which is, and which was, and which is to come,'' but this verse speaks of ``things'' that way, giving them an existence as eternal as God's, which is typical of the restoration.
				}

		\begin{framed}
		\scriptsize
			\begin{compactdesc}
				\item[JSP] and truth is knowledge of things as they are and as they were and as they are to come
				\item[RB2] and truth is knowledge of things as they are, and as they were, and as they are to come, 
				\item[RB1] And truth is knowledge of things as they are, and as they were, and as they are to come;
				\item[BC] ---
				\item[DC 1835] and truth is knowledge of things as they are, and as they were, and as they are to come;
			\end{compactdesc}
		\end{framed}

	% -----
	\subsubsection{Section 93:25} % *DC93-25 
	% -----
		\label{DC93-25}

		And whatsoever is more or less than this%
			\footnote{%
				% \label{}%
				See \hyperref[DC98-7]{DC 98:7}.
				}
		is the spirit%
			\footnote{%
				% \label{}%
				I think it's using ``spirit'' here in the sense we find it in \hyperref[ALMA-34-35]{Alma 34:35}.
				}
		of that wicked one who was a liar from the beginning.

		\begin{framed}
		\scriptsize
			\begin{compactdesc}
				\item[JSP] and whatsoever is more or less than these is the spirit of that wicked one who was a liar from the begining
				\item[BC] ---
				\item[RB2] and whatsoever is more or less than these is the spirit of that wicked one who was a liar from the beginning. 
				\item[RB1] and whatsoever is more or less than these is the spirit of that wicked one who was a liar from the beginning.
				% \item[DC 1835] 
				% \item[DC 1844] 
				% \item[]
			\end{compactdesc}
		\end{framed}

	% -----
	\subsubsection{Section 93:26} % *DC93-26 
	% -----
		\label{DC93-26}

		The Spirit of truth is of God. I am the Spirit of truth, and John bore record of me, saying: He received a fulness of truth, yea, even of all truth;

		\begin{framed}
		\scriptsize
			\begin{compactdesc}
				\item[JSP] the spirit of truth is of God, I am the spirit of truth, and John bear reccord of me say<ing> he received a fullness of truth yea even all truth
				\item[BC] ---
				\item[RB2] The spirit of truth is of God. I am the spirit of truth. and John bear reccord of me saying, he received a fullness of truth, yea even all truth 
				\item[RB1] [paragraph break] The Spirit of truth is of God: I am the spirit of truth: And John bore record of me, saying he received a fulness of truth; yea, even all truth. 
				% \item[DC 1835] 
				% \item[DC 1844] 
				% \item[]
			\end{compactdesc}
		\end{framed}

	% -----
	\subsubsection{Section 93:27} % *DC93-27 
	% -----
		\label{DC93-27}

		And no man receiveth a fulness unless he keepeth his commandments.

		\begin{framed}
		\scriptsize
			\begin{compactdesc}
				\item[JSP] and no man receiveth a fulness unless he keepeth his commandments
				\item[BC] ---
				\item[RB2] and no man receiveth a fullness unless he keepeth his commandments 
				\item[RB1] And no man receiveth a fulness unless he keepeth his commandments. 
				% \item[DC 1835] 
				% \item[DC 1844] 
				% \item[]
			\end{compactdesc}
		\end{framed}

	% -----
	\subsubsection{Section 93:28} % *DC93-28 
	% -----
		\label{DC93-28}

		He that keepeth his commandments receiveth truth and light, until he is glorified in truth and knoweth all things.

		\begin{framed}
		\scriptsize
			\begin{compactdesc}
				\item[JSP] he that keepeth his commandments receiveth truth and light untill he is glor[i]fied in truth and knoweth all things,
				\item[BC] ---
				\item[RB2] he that keepeth his commandments receiveth truth and light until he is glorified in truth and knoweth all things. 
				\item[RB1] He that keepeth his commandments, receiveth truth \& light, untill he is glorified in truth, and knoweth all things. 
				% \item[DC 1835] 
				% \item[DC 1844] 
				% \item[]
			\end{compactdesc}
		\end{framed}

	% -----
	\subsubsection{Section 93:29} % *DC93-29 
	% -----
		\label{DC93-29}

		Man was also in the beginning with God. Intelligence, or the light of truth, was not created or made, neither indeed can be.

		\begin{framed}
		\scriptsize
			\begin{compactdesc}
				\item[JSP] man was also in the begining with God, inteligence or the Light of truth was not created or made neith[er] indeed can be
				\item[BC] ---
				\item[RB2] Man was also in the beginning with God. Intelligence or the light of truth was not created or made neither indeed can be. 
				\item[RB1] Man was also in the beginning with God. Intelligence, or the light of truth, was not created or made, neither indeed can be. 
				% \item[DC 1835] 
				% \item[DC 1844] 
				% \item[]
			\end{compactdesc}
		\end{framed}

	% -----
	\subsubsection{Section 93:30} % *DC93-30 
	% -----
		\label{DC93-30}

		All truth is independent in that sphere in which God has placed it, to act for itself, as all intelligence also; otherwise there is no existence.

		\begin{framed}
		\scriptsize
			\begin{compactdesc}
				\item[JSP] all truth is independent in that sphere <in which God has placed it---> to act for itself as all inteligenc[e] also otherwise there is no existance
				\item[BC] ---
				\item[RB2] All truth is independent in that sphere in which God has placed it to act for itself. as all intelligence also otherwise there is no existence, 
				\item[RB1] All truth is independent in that sphere in which God has placed it to act for itself as all intelligence also, otherwise there is no existance.
				% \item[DC 1835] 
				% \item[DC 1844] 
				% \item[]
			\end{compactdesc}
		\end{framed}

	% -----
	\subsubsection{Section 93:31} % *DC93-31 
	% -----
		\label{DC93-31}

		Behold, here is the \hyperref[AGENCY]{agency} of man, and here is the condemnation of man; because that which was from the beginning is plainly manifest unto them, and they receive not the light.%
			\footnote{%
				% \label{}%
				Compare \hyperref[DC45-29]{DC 45:29}.
				}

		\begin{framed}
		\scriptsize
			\begin{compactdesc}
				\item[JSP] behold here is the agency of man and here is the condemnation of man because that which was from the begining is p[l]ainly manifest unto them and they receive not the light
				\item[BC] ---
				\item[RB2] behold here is the agency of man, and here is the condemnation of man, because that which was from the beginning is plainly manifest unto them, and they receive not the light, 
				\item[RB1] Behold here is the agency of man, and here is the condemnation of man: Because, that which was from the beginning, \sout{was} <is> plainly manifest unto them, and they receive not the light. 
				% \item[DC 1835] 
				% \item[DC 1844] 
				% \item[]
			\end{compactdesc}
		\end{framed}

	% -----
	\subsubsection{Section 93:32} % *DC93-32 
	% -----
		\label{DC93-32}

		And every man whose spirit receiveth not the light is under condemnation.

		\begin{framed}
		\scriptsize
			\begin{compactdesc}
				\item[JSP] and evry man whose spirit rec[e]iveth not the light
				\item[BC] ---
				\item[RB2] and every man whose spirit receiveth not the light, 
				\item[RB1] And every man whose Spirit receiveth not the light is under condemnation, 
				% \item[DC 1835] 
				% \item[DC 1844] 
				% \item[]
			\end{compactdesc}
		\end{framed}

	% -----
	\subsubsection{Section 93:33} % *DC93-33 
	% -----
		\label{DC93-33}

		For man is spirit.%
			\footnote{%
				% \label{}%
				Seems to connect to \hyperref[JOHN-4-24]{John 4:24}. We are a spiritual being---whatever that means, there's a distinction drawn here between it and ``element.''
				} 
		The elements are eternal, and spirit and element, inseparably connected, receive a fulness of joy;

		\begin{framed}
		\scriptsize
			\begin{compactdesc}
				\item[JSP] for man is spirit the Elements are eternal and spirit and element inseperably connected receiveth a fulness of Joy
				\item[BC] ---
				\item[RB2] for man is spirit, the Elements are eternal, and spirit and Element inseperably connected receiveth a fullness of Joy, 
				\item[RB1] for man is Spirit. The e$\diamond\diamond\diamond\diamond$ents are eternal: And spirit \& element inseperably connected \sout{connected} receiveth a fulness of Joy; 
				% \item[DC 1835] 
				% \item[DC 1844] 
				% \item[]
			\end{compactdesc}
		\end{framed}

	% -----
	\subsubsection{Section 93:34} % *DC93-34 
	% -----
		\label{DC93-34}

		And when separated, man cannot receive a fulness of joy.

		\begin{framed}
		\scriptsize
			\begin{compactdesc}
				\item[JSP] and when seperated man cannot receiv <a> fulness of Joy
				\item[BC] ---
				\item[RB2] and when separated man cannot receive a fullness of Joy.
				\item[RB1] and when seperated man cannot receive a fulness of Joy. 
				% \item[DC 1835] 
				% \item[DC 1844] 
				% \item[]
			\end{compactdesc}
		\end{framed}

	% -----
	\subsubsection{Section 93:35} % *DC93-35 
	% -----
		\label{DC93-35}

		The elements are the tabernacle of God; yea, man is the tabernacle of God, even temples; and whatsoever temple is defiled, God shall destroy that temple.%
			\footnote{%
				% \label{}%
				See this as an example of something in JS that looks like the idea of \hyperref[NATIVEELEMENT]{native element}?
				}

		\begin{framed}
		\scriptsize
			\begin{compactdesc}
				\item[JSP] the elements are the tabernacle of God, yea man is the tabernacle of God even temples and whatsoeve[r] temple is defiled God shall distroy that temple,
				\item[BC] ---
				\item[RB2] The Elements are the tabernacle of God. yea man is the tabernacle of God, even temples and whatsoever temple is defiled God shall destroy that temple 
				\item[RB1] The elements are the tabernacle of God; yea, man is the tabernacle of God even temples. And whatsoever temple is defiled, God shall destroy that temple. 
				% \item[DC 1835] 
				% \item[DC 1844] 
				% \item[]
			\end{compactdesc}
		\end{framed}

	% -----
	\subsubsection{Section 93:36} % *DC93-36 
	% -----
		\label{DC93-36}

		The glory of God is intelligence, or, in other words, light and truth.

		\begin{framed}
		\scriptsize
			\begin{compactdesc}
				\item[JSP] the glory of God is inteligence or in other words light \& truth
				\item[BC] ---
				\item[RB2] The glory of God is intelligence, or in other words, light and truth 
				\item[RB1] The glory of God is intelligence; or in other words, light, and \sout{truth} truth. 
				% \item[DC 1835] 
				% \item[DC 1844] 
				% \item[]
			\end{compactdesc}
		\end{framed}

	% -----
	\subsubsection{Section 93:37} % *DC93-37 
	% -----
		\label{DC93-37}

		Light and truth \hyperref[FORSAKE]{forsake} that evil one.

		\begin{framed}
		\scriptsize
			\begin{compactdesc}
				\item[JSP] light and truth forsaketh that evil one
				\item[BC] ---
				\item[RB2] light and truth forsaketh that evil one. 
				\item[RB1] Light <\& truth> forsaketh that evil one. 
				% \item[DC 1835] 
				% \item[DC 1844] 
				% \item[]
			\end{compactdesc}
		\end{framed}

	% -----
	\subsubsection{Section 93:38} % *DC93-38 
	% -----
		\label{DC93-38}

		Every spirit of man was innocent in the beginning; and God having redeemed man from the fall, men became again, in their infant state, innocent before God.

		\begin{framed}
		\scriptsize
			\begin{compactdesc}
				\item[JSP] evry spirit of man was innocent in the begining, and God having redeemed man from, the fall <man> became again in their infant state <innocent> before God
				\item[BC] ---
				\item[RB2] Every spirit of man was innocent in the beginning, and God haveing redeemed man from the fall, man became again in their infant state innocent before God, 
				\item[RB1] Every spirit of man was innocent in the beginning, And God having redeemed man from the fall, man became again in their infant state, innocent before God; 
				% \item[DC 1835] 
				% \item[DC 1844] 
				% \item[]
			\end{compactdesc}
		\end{framed}

	% -----
	\subsubsection{Section 93:39} % *DC93-39 
	% -----
		\label{DC93-39}

		And that wicked one
			\footnote{%
				% \label{}%
				Notice the similar language of ``coming'' and ``taking away'' in \hyperref[MARK-4-15]{Mark 4:15}, \hyperref[LUKE-8-12]{Luke 8:12}, 
				}%
		cometh and taketh away%
			\footnote{%
				% \label{}%
				Compare \hyperref[MATTHEW-13-19]{Matthew 13:19}, and notice the ways that this happens---there is an interaction between the seed, or what is planted, and the ``tradition'' of someone's fathers, or (broadly understood) what they already think. We have to try our best to integrate new stuff into the old and we have to be willing to give up things we thought were true when it looks like they aren't anymore. As Christ clarifies in that Matthew verse, whether people \textit{understand} the word (which is the \gr{λόγος}, by the way) makes a huge difference to the effect that it has on them---we have to keep this in mind when doing any kind of teaching. Compare \hyperref[DC68-25]{DC 68:25}.
				}
		light and truth, through disobedience,%
			\footnote{%
				% \label{}%
				See \hyperref[DC121-17]{DC 121:17}, 
				}
		from the children of men, and because%
			\footnote{%
				% \label{}%
				Notice that the ``traditions'' can also be good, of course---see \hyperref[ALMA-3-11]{Alma 3:11}, \hyperref[ALMA-9-8]{Alma 9:8}.
				}
		of the
			\footnote{%
				% \label{}%
				One thing that sticks out to me about certain passages involving these traditions---one is \hyperref[ALMA-21-16]{Alma 21:16--17}, but see the list below---is how much we need the Holy Ghost to be poured out upon a group of people in order to convince them that the traditions are false.
				}%
		tradition%
			\footnote{%
				% \label{}%
				There are many places in the scriptures where someone observes how damaging these ``traditions'' can be. See \hyperref[MATTHEW-16-12]{Matthew 16:12}, \hyperref[MARK-7-3]{Mark 7:2--23} (a focus on traditions can cause us to miss the principles that actually matter), \hyperref[GALATIANS-1-13]{Galatians 1:13--14}, \hyperref[JACOB-3-7]{Jacob 3:7--9}, \hyperref[MOSIAH-1-3]{Mosiah 1:3--5}, \hyperref[MOSIAH-10-12]{Mosiah 10:12--18}, \hyperref[MOSIAH-28-2]{Mosiah 28:2}, \hyperref[ALMA-3-8]{Alma 3:8--11}, \hyperref[ALMA-9-14]{Alma 9:14--17}, \hyperref[ALMA-17-9]{Alma 17:9}, \hyperref[ALMA-17-15]{Alma 17:15}, \hyperref[ALMA-18-2]{Alma 18:2--5} (notice that the traditions can be used in a constructive way to cause belief to go in a different direction), \hyperref[ALMA-19-14]{Alma 19:14}, \hyperref[ALMA-20-10]{Alma 20:10, 13}, \hyperref[ALMA-21-16]{Alma 21:16--17}, \hyperref[ALMA-23-3]{Alma 23:3, 5}; \hyperref[ALMA-24-7]{Alma 24:7--9}, \hyperref[ALMA-24-29]{Alma 24:29--30}, \hyperref[ALMA-25-6]{Alma 25:6}, \hyperref[ALMA-26-23]{Alma 26:23--24}, \hyperref[ALMA-30-14]{Alma 30:14--16}, \hyperref[ALMA-30-14]{Alma 30:23--31}, \hyperref[ALMA-31-16]{Alma 31:16--17, 22}; \hyperref[ALMA-37-9]{Alma 37:9}, \hyperref[ALMA-37-35]{Alma 37:35}, \hyperref[ALMA-60-32]{Alma 60:32--33}, \hyperref[HELAMAN-5-19]{Helaman 5:19}, \hyperref[HELAMAN-15-4]{Helaman 15:4--17}, \hyperref[3NEPHI-1-11]{3 Nephi 1:11} (this is a positive mention of the traditions), \hyperref[4NEPHI-1-38]{4 Nephi 1:38--39}, \hyperref[DC3-18]{DC 3:18}, \hyperref[DC74]{DC 74}, 
				}
		of their fathers.

		\begin{framed}
		\scriptsize
			\begin{compactdesc}
				\item[JSP] and that wicked one cometh and taketh away light and truth through disobeidienc [disobedience] from the children of men and becam[e] of the tradition of their fathers
				\item[BC] ---
				\item[RB2] and that wicked one cometh and taketh away light and truth through disobedience from the children of men, and because of the tradition of their fathers. 
				\item[RB1] and that wicked one cometh and taketh away light and truth through disobedience, from the children of men, and because of the tradition of their father's. 
				% \item[DC 1835] 
				% \item[DC 1844] 
				% \item[]
			\end{compactdesc}
		\end{framed}

	% -----
	\subsubsection{Section 93:40} % *DC93-40 
	% -----
		\label{DC93-40}

		But I have commanded you to bring up your children in light and truth.

		\begin{framed}
		\scriptsize
			\begin{compactdesc}
				\item[JSP] but I have commanded you to bring up your Children in light and truth,
				\item[BC] ---
				\item[RB2] But I have commanded you to bring up your children in light and truth, 
				\item[RB1] But I have commanded you to bring up your children in light \& truth: 
				% \item[DC 1835] 
				% \item[DC 1844] 
				% \item[]
			\end{compactdesc}
		\end{framed}

	% -----
	\subsubsection{Section 93:41} % *DC93-41 
	% -----
		\label{DC93-41}

		But verily I say unto you, my servant Frederick G. Williams, you have continued under this condemnation;

		\begin{framed}
		\scriptsize
			\begin{compactdesc}
				\item[JSP] but verily I say unto you my servant Frederick [G. Williams] you have continued under this condemnation
				\item[BC] ---
				\item[RB2] but verily I say unto you, my servant Frederick [G. Williams] you have continued under this condemnation 
				\item[RB1] But verily I say unto you my servant Frederick [G. Williams], you have continued under this condemnation; 
				% \item[DC 1835] 
				% \item[DC 1844] 
				% \item[]
			\end{compactdesc}
		\end{framed}

	% -----
	\subsubsection{Section 93:42} % *DC93-42 
	% -----
		\label{DC93-42}

		You have not taught your children light and truth, according to the commandments; and that wicked one hath power, as yet, over you, and this is the cause of your affliction.

		\begin{framed}
		\scriptsize
			\begin{compactdesc}
				\item[JSP] you have not taught your Children light and truth according to the Commandments and that wicked one hath power as yet over you and this is the caus of your affliction
				\item[BC] ---
				\item[RB2] you have not taught your children light and truth according to the commandments and that wicked one hath power as yet over you, and this is the cause of your affliction. 
				\item[RB1] you have not taught your children light and truth according to the commandments; and that wicked one hath power as yet over you; and this is the cause of your affliction. 
				% \item[DC 1835] 
				% \item[DC 1844] 
				% \item[]
			\end{compactdesc}
		\end{framed}

	% -----
	\subsubsection{Section 93:43} % *DC93-43 
	% -----
		\label{DC93-43}

		And now a commandment I give unto you---if you will be delivered you shall set in order your own house, for there are many things that are not right in your house.

		\begin{framed}
		\scriptsize
			\begin{compactdesc}
				\item[JSP] and now a commandment I give unto you and if ye will be delivered you shall set in order your own house for there are many things that are not right in your house
				\item[BC] ---
				\item[RB2] And now a commandment I give unto you, and if ye will be delivered you shall set in order your own house for there are many things that are not right in your house. 
				\item[RB1] And now, a commandment I give unto you, and if ye will be delivered, you shall set in order your own house; for there are many things that are not right in your house. 
				% \item[DC 1835] 
				% \item[DC 1844] 
				% \item[]
			\end{compactdesc}
		\end{framed}

	% -----
	\subsubsection{Section 93:44} % *DC93-44 
	% -----
		\label{DC93-44}

		Verily, I say unto my servant Sidney Rigdon, that in some things he hath not kept the commandments concerning his children; therefore, first set in order thy house.

		\begin{framed}
		\scriptsize
			\begin{compactdesc}
				\item[JSP] verely I say unto my servant Sidney [Rigdon] that in some things he hath not kept the commandments concerning his children therefore firstly set in order thy house,
				\item[BC] ---
				\item[RB2] Verily I say unto my servant Sidney [Rigdon] that in some things he hath not kept the commandments concerning his children. Therefore firstly set in order thy house. 
				\item[RB1] Verily I say, unto you my servant Sidney [Rigdon], that in some things he hath not kept the commandments concerning his children. Therefore firstly set in order thy house. 
				% \item[DC 1835] 
				% \item[DC 1844] 
				% \item[]
			\end{compactdesc}
		\end{framed}

	% -----
	\subsubsection{Section 93:45} % *DC93-45 
	% -----
		\label{DC93-45}

		Verily, I say unto my servant Joseph Smith, Jun., or in other words, I will call you friends, for you are my friends, and ye shall have an inheritance with me---

		\begin{framed}
		\scriptsize
			\begin{compactdesc}
				\item[JSP] and verely I say unto my servant Joseph (or in othe[r] words I will call you friends) for ye are my friends) and ye shall have an inheritance with me
				\item[BC] ---
				\item[RB2] And verily I say unto you my servant Joseph, or in otherwords, I will call you friends for ye are my friends and ye shall have an inheritance with me. 
				\item[RB1] And verily, I say unto my servant Joseph, or in other words I will call you friends, for ye are my friends, and ye shall have an inheritance with me. 
				% \item[DC 1835] 
				% \item[DC 1844] 
				% \item[]
			\end{compactdesc}
		\end{framed}

	% -----
	\subsubsection{Section 93:46} % *DC93-46 
	% -----
		\label{DC93-46}

		I called you servants for the world's sake, and ye are their servants for my sake---

		\begin{framed}
		\scriptsize
			\begin{compactdesc}
				\item[JSP] I called you servants for the worlds sake and ye are their servants for my sake
				\item[BC] ---
				\item[RB2] I called you servants for the worlds sake. and ye are their servants for my sake, 
				\item[RB1] I called you servants for the world's sake, and ye are their servants for my sake. 
				% \item[DC 1835] 
				% \item[DC 1844] 
				% \item[]
			\end{compactdesc}
		\end{framed}

	% -----
	\subsubsection{Section 93:47} % *DC93-47 
	% -----
		\label{DC93-47}

		And now, verily I say unto Joseph Smith, Jun.---You have not kept the commandments, and must needs stand rebuked before the Lord;

		\begin{framed}
		\scriptsize
			\begin{compactdesc}
				\item[JSP] and now verely I say unto you Joseph you have not kept the commandments and must needs stand rebuked before the lord
				\item[BC] ---
				\item[RB2] and now verily I say unto <you> Joseph you have not kept the commandments and must needs stand rebuked before the lord. 
				\item[RB1] And now verily I say unto you, Joseph, you have not kept the commandments, and must needs stand rebuked before the Lord. 
				% \item[DC 1835] 
				% \item[DC 1844] 
				% \item[]
			\end{compactdesc}
		\end{framed}

	% -----
	\subsubsection{Section 93:48} % *DC93-48 
	% -----
		\label{DC93-48}

		Your family must needs repent and forsake some things, and give more earnest heed unto your sayings, or be removed out of their place.

		\begin{framed}
		\scriptsize
			\begin{compactdesc}
				\item[JSP] your family must needs repent and forsake some things and give more earnest heed unto your sayings or be removed out of their place
				\item[BC] ---
				\item[RB2] Your family must needs repent and forsake some things. and give more earnest heed unto your sayings or be removed out of their place. 
				\item[RB1] Your family must needs repent and forsake some things and give more earnest heed unto your sayings, or be removed out of their place: 
				% \item[DC 1835] 
				% \item[DC 1844] 
				% \item[]
			\end{compactdesc}
		\end{framed}

	% -----
	\subsubsection{Section 93:49} % *DC93-49 
	% -----
		\label{DC93-49}

		What I say unto one I say unto all; pray always lest that wicked one have power in you, and remove you out of your place.

		\begin{framed}
		\scriptsize
			\begin{compactdesc}
				\item[JSP] what I say unto one I say unto all pray always lest that wicked one have power in you and remove you out of your place
				\item[BC] ---
				\item[RB2] what I say unto one I say unto all. pray always lest that wicked one have power \sout{over} <in> you, and remove you out of your place. 
				\item[RB1] What I say unto one, I say to all. Pray always, lets [lest] that wicked one have power in you, and remove you out of your place. 
				% \item[DC 1835] 
				% \item[DC 1844] 
				% \item[]
			\end{compactdesc}
		\end{framed}

	% -----
	\subsubsection{Section 93:50} % *DC93-50 
	% -----
		\label{DC93-50}

		My servant Newel K. Whitney also, a bishop of my church, hath need to be chastened, and set in order his family, and see that they are more diligent and concerned at home, and pray always, or they shall be removed out of their place.

		\begin{framed}
		\scriptsize
			\begin{compactdesc}
				\item[JSP] my servant Newel [K. Whitney] also the Bishop of my Church hath need to be chastened and set in order his family and see that they are more diligent and concerned at home and pray always or they shall be removed out of their place
				\item[BC] ---
				\item[RB2] My servant Newell [Newel K. Whitney] also. the Bishop of my church hath need to be chastened, and set in order his family, and see that they are more dilligent and concerned at home, and pray always or they shall be removed out of their place. 
				\item[RB1] My servant Newel [K. Whitney], also, the Bishop of my church hath need to be chastened, and set in order his family, and see that they are more dilligent and concerned at home, and pray always, or they shall be removed out of their place. 
				% \item[DC 1835] 
				% \item[DC 1844] 
				% \item[]
			\end{compactdesc}
		\end{framed}

	% -----
	\subsubsection{Section 93:51} % *DC93-51 
	% -----
		\label{DC93-51}

		Now, I say unto you, my friends, let my servant Sidney Rigdon go on his journey, and make haste, and also proclaim the acceptable year of the Lord, and the gospel of salvation, as I shall give him utterance; and by your prayer of faith with one consent I will uphold him.

		\begin{framed}
		\scriptsize
			\begin{compactdesc}
				\item[JSP] now I say unto you my friends let my servant Sidny go his Journey and make haste and also proclaim the acceptable year of the Lord and the gospel of salvation as I shall give him utterence and by your prayr of faith with one consent I will uphold him
				\item[BC] ---
				\item[RB2] Now I say unto you my friends, let my servant Sidney go his journey and make haste and also proclaim the acceptable year of the Lord and the gospel of salvation as I shall give him utterance and by your prayer of faith with one consent. I will uphold him. 
				\item[RB1] Now I say unto you my friends, let my servant Sidney go his Journey, and make haste, and also, proclaim the acceptable year of the Lord, and the gospel of Salvation, as I shall give him utterance. And by your prayer of faith, with one consent, I will uphold him. 
				% \item[DC 1835] 
				% \item[DC 1844] 
				% \item[]
			\end{compactdesc}
		\end{framed}

	% -----
	\subsubsection{Section 93:52} % *DC93-52 
	% -----
		\label{DC93-52}

		And let my servants Joseph Smith, Jun., and Frederick G. Williams make haste also, and it shall be given them even according to the prayer of faith; and inasmuch as you keep my sayings you shall not be confounded in this world, nor in the world to come.

		\begin{framed}
		\scriptsize
			\begin{compactdesc}
				\item[JSP] and let my servants Joseph \& Frederick make haste also and it shall be given them even according to the prayer of faith [paragraph break] and inasmuch as you keep my sayings you shall not be confounded in this world nor in the world to come \sout{Am}
				\item[BC] ---
				\item[RB2] and let my servants Joseph and Frederick make haste also, and it shall be given them even according to the prayer of faith. And in as much as you keep my sayings you shall not be confounded in this world nor in the world to come--- 
				\item[RB1] And let my Servants Joseph and Frederick make haste also, and it shall be given them even according to the prayer of faith. And in as much as you keep my sayings, you shall not be confounded, in this world, nor in the world to come.
				% \item[DC 1835] 
				% \item[DC 1844] 
				% \item[]
			\end{compactdesc}
		\end{framed}

	% -----
	\subsubsection{Section 93:53} % *DC93-53 
	% -----
		\label{DC93-53}

		And, verily I say unto you, that it is my will that you should hasten to translate my scriptures, and to obtain a knowledge of history, and of countries, and of kingdoms, of laws of God and man, and all this for the salvation of Zion. Amen.

		\begin{framed}
		\scriptsize
			\begin{compactdesc}
				\item[JSP] and verely I say unto you that it is my will that ye should hasten to translate my Scriptures and to obtain a knowledg of history and of Countries and of Kingdoms and of laws, of \sout{man} \sout{\& of God} <God \& man> and all these for the salvation of Zion Amen
				\item[BC] ---
				\item[RB2] And verily I say unto you that it is my will that ye should hasten to translate my scriptures and to obtain a knowledge of history and of countries and of Kingdoms and of laws of God and \sout{men} <man>, and all this for the salvation of Zion. Amen
				\item[RB1] And verily I say unto you, that it is my will that ye should hasten to translate my scriptures, and to obtain a knowledge of history, and of countries, \& of kingdoms and of laws of God \& of man; and all this for the Salvation of Zion: Amen. 
				% \item[DC 1835] 
				% \item[DC 1844] 
				% \item[]
			\end{compactdesc}
		\end{framed}